\section{Overview}

In recent years, the field of manufacturing has witnessed significant advancements, particularly with the rise of additive manufacturing techniques such as 3D printing. This technology has revolutionized the way products are designed and produced, enabling rapid prototyping, customization, and the creation of complex geometries that were previously unattainable with traditional manufacturing methods. The ability to quickly turn digital designs into physical models has opened up new possibilities for innovation across various industries, from aerospace to healthcare. However, despite its advantages, 3D printing is not without its limitations. It can be slower than traditional manufacturing for large-scale production, and some materials used in 3D printing may not offer the same strength or durability as those produced through conventional techniques.

This project focuses on the integrated design of an ESP32 TTGO T-Display device with a custom 3D-printed housing. The ESP32 TTGO T-Display is a versatile microcontroller board that features built-in Wi-Fi and Bluetooth capabilities, making it ideal for a wide range of applications, including IoT devices, wearable technology, and interactive displays. By designing a custom housing for the ESP32 TTGO T-Display using 3D printing, we aim to enhance the functionality and aesthetics of the device while also exploring the design for assembly principles to ensure ease of manufacturing and maintenance. The project will involve the entire design process, from conceptualization and CAD modeling to prototyping and testing, with a focus on optimizing the design for both performance and manufacturability. Through this project, we hope to demonstrate the potential of integrating 3D printing with electronic device design to create innovative and functional products.